% Part: computability
% Chapter: recursive-functions
% Section: partial-functions

\documentclass[../../../include/open-logic-section]{subfiles}

\begin{document}

\olfileid[tr]{cmp}{rec}{par}
\olsection{Kısmi Özyinelemeli Fonksiyonlar}

Özyinelemeli fonksiyonların tanımını gerekçelendirmek için, ilkel
özyinelemeli olmayan hesaplanabilir fonksiyonların bulunduğuna ilişkin
ispatımızın aslında çok daha fazlasını gösterdiğine dikkat edin. Sav basitti:
Kullandığımız tek şey, $f_x(y)$ ifadesi $x$ ile $y$'nin bir fonksiyonu olarak
hesaplanabilir olacak biçimde $f_0,f_1,\dots$ fonksiyonlarını
numaralandırabilmekti. Dolayısıyla sav, \emph{bu biçimde numaralandırılabilen
her fonksiyon sınıfına} uygulanır. Bu bizi açmaza sokar: Hesaplanabilir
fonksiyonları açıkça betimlemek istiyoruz, fakat hesaplanabilir
fonksiyonlardan oluşan bir topluluğun hiçbir açık betimi tüketici olamaz.

Çıkış yolu \emph{kısmi} fonksiyonları hesaba katmaktır. Kısmi hesaplanabilir
fonksiyonların numaralandırılmasının \emph{mümkün} olduğunu göreceğiz.
\iftag{TMs}{ Aslında Turing makinelerini sistematik biçimde
numaralandırabildiğimizden bunun böyle olduğunu neredeyse şimdiden
biliyoruz.}{} Daha sonra köşegen savımıza dönecek ve kısmi fonksiyonlar
katıldığında neden işlemediğini inceleyeceğiz.

Şimdi soru şudur: Bütün kısmi özyinelemeli fonksiyonları elde etmek için
ilkel özyinelemeli fonksiyonlara ne eklemeliyiz? İki şey yapmamız gerekir:
\begin{enumerate}
\item İlkel özyinelemeli fonksiyonlar tanımımızı kısmi fonksiyonlara da izin
  verecek biçimde değiştirmek.
\item Tanıma, bazı yeni kısmi fonksiyonları kapsayacak bir şey
  \emph{eklemek}.
\end{enumerate}

İlki kolaydır. Önceki gibi sıfır, ardıl ve izdüşümlerle başlayacak; bileşke
ile ilkel özyineleme altında kapanışı alacağız. Tek fark, tanımdaki bazı
terimlerin tanımsız olabilmesine izin verecek biçimde bileşke ve ilkel
özyineleme tanımlarını değiştirmemizdir. $f$ ile $g$ kısmi fonksiyonlarsa,
$f$'nin $x$'te tanımlı, yani $x$'in $f$'nin tanım kümesinde olduğunu
$f(x)\fdefined$ ile; tersini, yani $f$'nin $x$'te tanımsız olduğunu
$f(x)\fundefined$ ile göstereceğiz. $f(x)\simeq g(x)$ gösterimi, ya iki
ifadenin de tanımsız ya da ikisinin de tanımlı ve eşit olduğunu söyler. Bu
gösterimleri daha karmaşık terimler için de kullanacağız. $h$ ile
$g_0,\dots,g_k$ kısmi fonksiyonlarsa
\[
  h(g_0(\vec x),\dots,g_k(\vec x))
\]
ifadesinin ancak ve ancak her $g_i$, $\vec x$ noktasında ve $h$ de
$g_0(\vec x),\dots,g_k(\vec x)$ değerlerinde tanımlıysa tanımlı olduğu
uzlaşımını benimseyeceğiz. Bu anlayışla kısmi fonksiyonlar için bileşke ve
ilkel özyineleme tanımları yukarıdakilerle aynıdır; yalnızca ``$=$'' yerine
``$\simeq$'' koymamız gerekir.

Kısmi fonksiyonları elde etmek üzere ilkel özyinelemeli fonksiyonlar tanımına
ekleyeceğimiz şey \emph{sınırsız arama işlecidir}. Doğal sayılar üzerinde
herhangi bir kısmi $f(x,\vec z)$ fonksiyonu için $\mu x\;f(x,\vec z)$
ifadesini, böyle bir $x$ varsa
\begin{quote}
  $f(0,\vec z),f(1,\vec z),\dots,f(x,\vec z)$ değerlerinin tümünün tanımlı
  ve $f(x,\vec z)=0$ olduğu en küçük $x$
\end{quote}
olarak tanımlayalım; aksi durumda $\mu x\;f(x,\vec z)$ tanımsızdır. Bu,
$\mu x\;f(x,\vec z)$ ifadesini tek türlü tanımlar.

\begin{explain}
Tanımımızın \iftag{TMs}{Turing makinelerine,}{} algoritmalara ya da belirli
bir hesaplama modeline başvurmadığına dikkat edin. Yine de bileşke ve ilkel
özyinelemede olduğu gibi sınırsız aramanın arkasında işlemsel, hesaplamalı bir
sezgi vardır. Kısmi bir fonksiyonun hesaplanabilirliği söz konusu olduğunda,
fonksiyonun tanımsız olduğu argümanlar hesabın durmadığı girdilere karşılık
gelir. $\mu x\;f(x,\vec z)$ değerini hesaplama yordamı şöyledir:
$0$ döndüren bir değer bulununcaya kadar
$f(0,\vec z),f(1,\vec z),f(2,\vec z),\dots$ hesaplanır. Fakat aradaki
hesaplardan herhangi biri durmazsa $\mu x\;f(x,\vec z)$ hesabı da durmaz.
\end{explain}

$R(x,\vec z)$ herhangi bir bağıntıysa $\mu x\;R(x,\vec z)$ ifadesini
$\mu x\;(1\tsub\Char{R}(x,\vec z))$ olarak tanımlarız. Başka bir deyişle
$\mu x\;R(x,\vec z)$, $R(x,\vec z)$ koşulunu sağlayan en küçük $x$ değerini
döndürür. Dolayısıyla $f(x,\vec z)$ tümel bir fonksiyonsa
$\mu x\;f(x,\vec z)$ ile $\mu x\;(f(x,\vec z)=0)$ aynıdır. Fakat asıl
tanımımızın daha genel olduğuna dikkat edin; çünkü $f(x,\vec z)$ fonksiyonunun
her yerde tanımlı olmamasına izin verir. Buna karşılık bir bağıntının
karakteristik fonksiyonu her zaman tümeldir.

\begin{defn}
\emph{Kısmi özyinelemeli fonksiyonlar} kümesi, doğal sayıların çeşitli
kuvvetlerinden doğal sayılara giden kısmi fonksiyonlar arasında sıfır, ardıl
ve izdüşümleri içeren ve bileşke, ilkel özyineleme ile sınırsız arama altında
kapalı olan en küçük kümedir.
\end{defn}

Elbette kısmi özyinelemeli fonksiyonların bazıları tümel, yani her argümanda
tanımlı olacaktır.

\begin{defn}
\ollabel{defn:recursive-fn}
\emph{Özyinelemeli fonksiyonlar} kümesi, tümel olan kısmi özyinelemeli
fonksiyonlar kümesidir.
\end{defn}

Bir özyinelemeli fonksiyonun her yerde tanımlı olduğunu vurgulamak için ona
bazen ``tümel özyinelemeli'' denir.

\end{document}

